% ------------------------------------------------------------
% FDI setup
% ------------------------------------------------------------

\pagestyle{empty}
\setlength{\parindent}{0pt}

\newcounter{fdipage}

% ------------------------------------------------------------
% Comic folder and blank comic pages
% ------------------------------------------------------------
% Enthaelt diese Datei und das Blanko-PDF, also alles, was fuer
% alle Sprachen gleich ist. Bei einer neuen Blanko-Version nur
% \FDIBlankoFile anpassen.

\newcommand{\FDIComicFolder}{comic_files}
\newcommand{\FDIBlankoFile}{\FDIComicFolder/BLANKO_1-12_v1.0.0.pdf}

% ------------------------------------------------------------
% Fonts setup depending on language
% ------------------------------------------------------------

% Normale Fonts
\newfontfamily{\FDIKarolingischeMinuskelDefault}{Karolingische_Minuskel.ttf}[
  Path = fonts/,
  FakeBold = 2
]

\newfontfamily{\FDIOverlayFontDefault}{KOMTXT__.ttf}[
  Path = fonts/
]

\newfontfamily{\FDIArialBOLDFontDefault}{ArialCEMTBlack.ttf}[
  Path = fonts/
]

\newfontfamily{\FDIArialFontDefault}{ARIAL.ttf}[
  Path = fonts/
]

% Elbische FDI-Schriften
% Wichtig: Diese Fonts werden NICHT global auf das Dokument angewendet,
% sondern nur über die FDI-Font-Makros unten verwendet.
% Renderer=Basic ist bei alten Tengwar-/Symbolfonts oft wichtig.

% Normale FDI-Sprechblasen / Standard-Overlay
\newfontfamily{\FDIElbArialFont}{tngani.ttf}[
  Path = fonts/,
  Renderer = Basic,
  Scale = 0.78
]

% Fette FDI-Sprechblasen / Hervorhebungen
\newfontfamily{\FDIElbArialBOLDFont}{tngani.ttf}[
  Path = fonts/,
  Renderer = Basic,
  Scale = 0.78,
  FakeBold = 1.8
]

% FDI-Overlay-Schrift, etwas größer für Comic-Texte
\newfontfamily{\FDIElbOverlayFont}{tngani.ttf}[
  Path = fonts/,
  Renderer = Basic,
  Scale = 0.82
]

% FDI-Karolingische-Minuskel-Ersatzschrift
\newfontfamily{\FDIElbKarolingischeMinuskel}{tngani.ttf}[
  Path = fonts/,
  Renderer = Basic,
  Scale = 0.86,
  FakeBold = 1.3
]

% ------------------------------------------------------------
% Public font aliases
% ------------------------------------------------------------

\providecommand{\KarolingischeMinuskel}{}
\providecommand{\FDIOverlayFont}{}
\providecommand{\ArialBOLDFont}{}
\providecommand{\ArialFont}{}

% ------------------------------------------------------------
% Apply FDI font aliases depending on language
% ------------------------------------------------------------
% Wichtig:
% Dieser Schalter verändert NICHT \setmainfont, \setsansfont oder \setmonofont.
% Er betrifft ausschließlich die FDI-Font-Makros:
% \KarolingischeMinuskel, \FDIOverlayFont, \ArialBOLDFont und \ArialFont.

\newcommand{\FDISetupLanguageFonts}{%
  \ifdefstring{\FDILanguage}{elb}{%
    \typeout{FDI: Language is elb. Switching FDI font macros only.}%

    % Nur FDI-Font-Makros auf elbische Varianten setzen.
    % Keine globalen Dokumentfonts ändern!
    \let\KarolingischeMinuskel\FDIElbKarolingischeMinuskel
    \let\FDIOverlayFont\FDIElbOverlayFont
    \let\ArialBOLDFont\FDIElbArialBOLDFont
    \let\ArialFont\FDIElbArialFont
  }{%
    \typeout{FDI: Language is not elb. Using default FDI font macros.}%

    % Normale FDI-Font-Makros setzen.
    \let\KarolingischeMinuskel\FDIKarolingischeMinuskelDefault
    \let\FDIOverlayFont\FDIOverlayFontDefault
    \let\ArialBOLDFont\FDIArialBOLDFontDefault
    \let\ArialFont\FDIArialFontDefault
  }%
}

% ------------------------------------------------------------
% Language folder depending on language
% Usage result: language_files/de, language_files/en, ...
% Enthaelt die Sprachdatei, den zugehoerigen pages-Ordner und
% alle sprachabhaengigen Bilder.
% ------------------------------------------------------------

\newcommand{\FDILanguageFolder}{%
  language_files/\FDILanguage%
}

% ------------------------------------------------------------
% Input language text file
% Loads: language_files/de/de.tex, language_files/en/en.tex, ...
% ------------------------------------------------------------

\newcommand{\FDIInputLanguageFile}{%
  \IfFileExists{\FDILanguageFolder/\FDILanguage.tex}{%
    \input{\FDILanguageFolder/\FDILanguage.tex}%
  }{%
    \PackageError{FDI}{Language file '\FDILanguageFolder/\FDILanguage.tex' not found}{%
      Check \string\FDILanguage\space and create the corresponding file.%
    }%
  }%
}

% ------------------------------------------------------------
% Set page folder depending on language
% Usage result: language_files/de/pages_de, language_files/en/pages_en, ...
% ------------------------------------------------------------

\newcommand{\FDIPageFolder}{%
  \FDILanguageFolder/pages_\FDILanguage%
}

% ------------------------------------------------------------
% Input page by number
% Usage: \FDIInputPage{1}
% Loads: language_files/de/pages_de/page_1.tex, ...
% ------------------------------------------------------------

\newcommand{\FDIInputPage}[1]{%
  \IfFileExists{\FDIPageFolder/page_#1.tex}{%
    \input{\FDIPageFolder/page_#1.tex}%
  }{%
    \PackageError{FDI}{Page file '\FDIPageFolder/page_#1.tex' not found}{%
      Check \string\FDILanguage\space and create the corresponding page file.%
    }%
  }%
}

% ------------------------------------------------------------
% Whole comic in one language
% Usage: \comic{de}        alle Comicseiten auf Deutsch
%        \comic[3]{en}     nur die ersten 3 Seiten, also Seite 1 bis 3
%        \comic[3-5]{en}   Seite 3 bis 5
%        \comic[7-7]{en}   nur Seite 7
% ------------------------------------------------------------
% Gibt alle Comicseiten der Sprache #2 direkt hintereinander aus.
% Sprache, Fonts und Texte gelten nur innerhalb von \comic. Nach
% \comic gelten wieder die Einstellungen des Dokuments, also
% \FDILanguage sowie \title/\subtitle/\shorttitle aus der in der
% Praeambel geladenen Sprachdatei.
% Dadurch koennen auch mehrere Sprachen nacheinander ausgegeben
% werden, z.B.:
%   \comic{de}
%   \comic{en}

% Anzahl der Comicseiten pro Sprache
\newcommand{\FDIComicPageCount}{12}

\makeatletter

\newcounter{FDIcomicpage}

% Optionales Argument von \comic auswerten.
% Aufruf immer mit angehaengtem --\FDI@stop, dadurch gilt:
%   "5"    ->  #1=5  #2=leer  ->  Seite 1 bis 5
%   "3-5"  ->  #1=3  #2=5     ->  Seite 3 bis 5
\def\FDI@stop{}

\def\FDI@parserange#1-#2-#3\FDI@stop{%
  \def\FDI@tmp{#2}%
  \ifx\FDI@tmp\@empty
    \def\FDI@firstpage{1}%
    \edef\FDI@lastpage{#1}%
  \else
    \edef\FDI@firstpage{#1}%
    \edef\FDI@lastpage{#2}%
  \fi
}

\newcommand{\comic}[2][\FDIComicPageCount]{%
  \begingroup
    % \edef, damit auch \comic{\FDILanguage} funktioniert
    \edef\FDILanguage{#2}%
    \FDISetupLanguageFonts
    \FDIInputLanguageFile
    % Seitenbereich bestimmen
    \edef\FDI@tmp{#1}%
    \expandafter\FDI@parserange\FDI@tmp--\FDI@stop
    % Beide Zaehler auf die Seite vor der ersten gewuenschten Seite setzen.
    % fdipage waehlt die Seite im Hintergrund-PDF und laeuft deshalb mit,
    % damit \comic[3-5] auch die Zeichnungen der Seiten 3 bis 5 zeigt.
    \setcounter{FDIcomicpage}{\FDI@firstpage}%
    \addtocounter{FDIcomicpage}{-1}%
    \setcounter{fdipage}{\value{FDIcomicpage}}%
    \@whilenum\value{FDIcomicpage}<\FDI@lastpage\do{%
      \stepcounter{FDIcomicpage}%
      \expandafter\FDIInputPage\expandafter{\number\value{FDIcomicpage}}%
    }%
    % letzte Comicseite ausgeben, solange die Sprache noch gesetzt ist
    \clearpage
  \endgroup
}

\makeatother

% ------------------------------------------------------------
% FDI header/footer offsets
% ------------------------------------------------------------
% Positive X-Werte verschieben nach rechts.
% Negative X-Werte verschieben nach links.
% Positive Y-Werte verschieben nach oben.
% Negative Y-Werte verschieben nach unten.

\newlength{\FDIHeaderXOffset}
\newlength{\FDIHeaderYOffset}
\newlength{\FDIFooterXOffset}
\newlength{\FDIFooterYOffset}

\setlength{\FDIHeaderXOffset}{0pt}
\setlength{\FDIHeaderYOffset}{20pt}
\setlength{\FDIFooterXOffset}{0pt}
\setlength{\FDIFooterYOffset}{-35pt}

% ------------------------------------------------------------
% FDI pagestyle: original inggrid pagestyle with offsets
% ------------------------------------------------------------

\makeatletter

\let\FDI@original@ps@inggrid\ps@inggrid

\def\ps@inggridfdi{%
  \FDI@original@ps@inggrid%

  \let\FDI@oddhead\@oddhead%
  \let\FDI@evenhead\@evenhead%
  \let\FDI@oddfoot\@oddfoot%
  \let\FDI@evenfoot\@evenfoot%

  \def\@oddhead{%
    \hspace*{\FDIHeaderXOffset}%
    \raisebox{\FDIHeaderYOffset}[0pt][0pt]{\FDI@oddhead}%
    \hspace*{-\FDIHeaderXOffset}%
  }%

  \def\@evenhead{%
    \hspace*{\FDIHeaderXOffset}%
    \raisebox{\FDIHeaderYOffset}[0pt][0pt]{\FDI@evenhead}%
    \hspace*{-\FDIHeaderXOffset}%
  }%

  \def\@oddfoot{%
    \hspace*{\FDIFooterXOffset}%
    \raisebox{\FDIFooterYOffset}[0pt][0pt]{\FDI@oddfoot}%
    \hspace*{-\FDIFooterXOffset}%
  }%

  \def\@evenfoot{%
    \hspace*{\FDIFooterXOffset}%
    \raisebox{\FDIFooterYOffset}[0pt][0pt]{\FDI@evenfoot}%
    \hspace*{-\FDIFooterXOffset}%
  }%
}

\makeatother


% ------------------------------------------------------------
% FDI pages
% ------------------------------------------------------------

\newcommand{\FDIPage}[1]{%
  \clearpage
  \stepcounter{fdipage}%
  \thispagestyle{inggridfdi}%
  \AddToShipoutPictureBG*{%
    \begin{tikzpicture}[remember picture, overlay]
      \node[anchor=center, inner sep=0pt]
        at (current page.center) {%
          \includegraphics[
            page=\value{fdipage},
            width=\paperwidth,
            height=\paperheight
          ]{\FDIBlankoFile}%
        };
    \end{tikzpicture}%
  }%
  #1%
  \null
}

% ------------------------------------------------------------
% FDI Texts
% ------------------------------------------------------------
% Globale Variable für den Zeilenabstand
\newlength{\FDITextLineSkip}
\setlength{\FDITextLineSkip}{12pt}


% \FDIStyledText{schriftgröße}{ignored}{text}
% Achtung: kein äußerer {...}-Block, damit \baselineskip bis zum Absatzende wirkt.
\newcommand{\FDIStyledText}[3]{%
  \FDIOverlayFont
  \fontsize{#1}{\FDITextLineSkip}\selectfont
  \setlength{\baselineskip}{\FDITextLineSkip}%
  #3%
}

% \FDIText{x}{y}{breite}{text}
\newcommand{\FDIText}[4]{%
  \node[
    anchor=north west,
    inner sep=0pt,
    outer sep=0pt,
    minimum width=#3
  ] at ([xshift=#1,yshift=-#2]current page.north west) {%
    \begin{minipage}[t]{#3}
      \centering
      \setlength{\lineskip}{0pt}%
      \setlength{\parskip}{0pt}%
      \FDIOverlayFont
      #4%
    \end{minipage}%
  };
}

% \FDITextLeft{x}{y}{breite}{text}
\newcommand{\FDITextLeft}[4]{%
  \node[
    anchor=north west,
    inner sep=0pt,
    outer sep=0pt,
    minimum width=#3
  ] at ([xshift=#1,yshift=-#2]current page.north west) {%
    \begin{minipage}[t]{#3}
      \raggedright
      \setlength{\lineskip}{0pt}%
      \setlength{\parskip}{0pt}%
      \FDIOverlayFont
      #4%
    \end{minipage}%
  };
}

\newcommand{\FDIStyledTextMinuskel}[3]{%
  {\KarolingischeMinuskel\fontsize{#1}{#2}\selectfont #3}%
}

% \FDITextLS{x}{y}{breite}{schriftgröße}{zeilenabstand}{text}
% \FDIText{x}{y}{breite}{schriftgröße}{zeilenabstand}{text}
\newcommand{\FDITextLS}[6]{%
  \node[
    anchor=north west,
    text width=#3,
    align=center,
    font=\FDIOverlayFont,
    inner sep=0pt,
    outer sep=0pt,
    minimum width=#3,
    font=\FDIOverlayFont\fontsize{#4}{#5}\selectfont
  ] at ([xshift=#1,yshift=-#2]current page.north west) {#6};
}



% \FDIOvalWrapText{x}{y}{breite}{text}
\newcommand{\FDIOvalText}[4]{%
  % unsichtbare / randlose Sprechblasenfläche
  \node[
    ellipse,
    draw=none,
    minimum width=#3,
    minimum height=1.9cm,
    inner sep=0pt,
    outer sep=0pt
  ] at ([xshift=#1,yshift=-#2]current page.north west) {};

  % Text mit angenäherter ovaler Kontur
  \node[
    anchor=center,
    inner sep=0pt,
    outer sep=0pt
  ] at ([xshift=#1,yshift=-#2]current page.north west) {%
    \begin{minipage}{#3}
      \centering
      \FDIOverlayFont
      \parshape=5
        0.55cm \dimexpr#3-1.10cm\relax
        0.25cm \dimexpr#3-0.50cm\relax
        0.00cm #3
        0.25cm \dimexpr#3-0.50cm\relax
        0.55cm \dimexpr#3-1.10cm\relax
      #4%
    \end{minipage}%
  };
}

% \FDIArcTextOutside{x}{y}{start angle}{end angle}{radius}{text}{style}
% Text ist von außen lesbar.
\newcommand{\FDIArcTextOutside}[7]{%
  \path[
    decorate,
    decoration={
      text along path,
      text={|#7| #6},
      text align=center
    }
  ]
  ([xshift=#1,yshift=-#2]current page.north west)
  arc[start angle=#3,end angle=#4,radius=#5];
}

% \FDIArcTextInside{x}{y}{start angle}{end angle}{radius}{text}{style}
% Text ist von innen lesbar.
\newcommand{\FDIArcTextInside}[7]{%
  \path[
    decorate,
    decoration={
      text along path,
      reverse path,
      text={|#7| #6},
      text align=center
    }
  ]
  ([xshift=#1,yshift=-#2]current page.north west)
  arc[start angle=#3,end angle=#4,radius=#5];
}

% \FDITextRotated{x}{y}{breite}{winkel}{text}
% winkel: positiver Wert = gegen den Uhrzeigersinn, negativer Wert = im Uhrzeigersinn
\newcommand{\FDITextRotated}[5]{%
  \node[
    anchor=north west,
    text width=#3,
    align=center,
    inner sep=0pt,
    outer sep=0pt,
    rotate=#4
  ] at ([xshift=#1,yshift=-#2]current page.north west) {#5};
}

% ------------------------------------------------------------
% FDI Text Posistion Makros
% ------------------------------------------------------------
% \def statt \newcommand, damit dieselbe Sprachdatei mehrfach bzw.
% mehrere Sprachdateien nacheinander geladen werden koennen (siehe \comic).
\newcommand{\FDISetText}[4]{%
  \expandafter\def\csname s#1p#2t#3\endcsname{#4}%
}

\newcommand{\FDIUseText}[3]{%
  \csname s#1p#2t#3\endcsname
}

% ------------------------------------------------------------
% FDI city dots
% ------------------------------------------------------------
% \FDICityDot{x}{y}{größe}{füllfarbe}
% Setzt einen Punkt mit weißer Outline.
% Alle Werte werden vom oberen linken Seitenrand gemessen.
% Beispiel: \FDICityDot{7.2cm}{11.4cm}{7pt}{red!85!black}

\newlength{\FDICityDotOutlineWidth}
\setlength{\FDICityDotOutlineWidth}{1.2pt}

\tikzset{
  FDICityDotStyle/.style={
    circle,
    draw=white,
    line width=\FDICityDotOutlineWidth,
    inner sep=0pt,
    outer sep=0pt
  }
}

\newcommand{\FDICityDot}[4]{%
  \node[
    FDICityDotStyle,
    minimum size=#3,
    fill=#4
  ] at ([xshift=#1,yshift=-#2]current page.north west) {};
}


% ------------------------------------------------------------
% FDI movable line numbers / panel numbers
% Uses original inggrid / lineno counter
% ------------------------------------------------------------

\makeatletter

% Zentraler Stil der Nummern aus lineno / inggrid
\newcommand{\FDILineNumberStyle}{%
  \linenumberfont
}

% \FDILineNumberBox{x}{y}
% x = Abstand vom linken Seitenrand
% y = Abstand vom oberen Seitenrand
% Nutzt den originalen linenumber-Zähler fortlaufend weiter.
\newcommand{\FDILineNumberBox}[2]{%
  \stepcounter{linenumber}%
  \node[
    anchor=north west,
    inner sep=0pt,
    outer sep=0pt
  ] at ([xshift=#1,yshift=-#2]current page.north west) {%
    {\FDILineNumberStyle\thelinenumber}%
  };%
}

\makeatother


% ------------------------------------------------------------
% Horizontal lines
% ------------------------------------------------------------

\newcommand{\FDIHLine}[3]{%
  \draw[lightgray, line width=0.4pt]
    ([xshift=#1,yshift=-#3]current page.north west)
    --
    ([xshift=#2,yshift=-#3]current page.north west);
}


% ------------------------------------------------------------
% Point to point lines
% ------------------------------------------------------------

% \FDILine{x1}{y1}{x2}{y2}
% zieht eine Linie von Punkt 1 zu Punkt 2
% alle Werte werden von oben links gemessen
\newcommand{\FDILine}[4]{%
  \draw[black, line width=0.4pt]
    ([xshift=#1,yshift=-#2]current page.north west)
    --
    ([xshift=#3,yshift=-#4]current page.north west);
}